\scanpage{171}
% archive foliation: 170
\uline{Groupe fondamental} : (\unsure{Revêtts} infinis)

$S$ préschéma. Pour tout groupe discret $G$, on définit
\[
H^1_C(S,G)=\varinjlim_{T/S\in C}H^1(T/S,G)
\]
\deleted{$T$ \unsure{fid plat de type f. sur $S$}}\quad
\deleted{\unsure{fid plat sur $S$}}

où
\[
H^1(T/S,G)=H^1(K_{T/S},G)=H^1(\pi_0(K_{T/S}),G)
\]
où $C$ est une catégorie de $S$-préschémas.
Les cas les plus importants sont les suivants
\diagram{\[
\begin{array}{l}
\left\{\begin{array}{lll}
\text{(i)}&C_1&\text{$S$-préschémas fid plats quasi-compacts, \deleted{\unsure{loc. noeth.}}}\\[3pt]
\multicolumn{3}{l}{\left.\begin{array}{lll}
\text{(ii)}&C_2&\rule{3cm}{.3pt}\quad\text{de type fini}\\[3pt]
\text{(iii)}&C_3&\rule{3cm}{.3pt}\quad\text{quasi-finis}
\end{array}\right\}}
\end{array}\right.\\[6pt]
\left.\begin{array}{lll}
\text{(iv)}&C_4&\rule{3cm}{.3pt}\quad\text{finis par morceaux}\\[3pt]
\text{(v)}&C_5&\rule{3cm}{.3pt}\ \deleted{\unsure{étales}}\quad\text{finis étales p. morceaux}
\end{array}\right\}
\end{array}
\]}

\margin{élucider ici les questions d'effectivité des données de descente\ldots\ \unsure{les ramener} à la descente des morphismes étales de type fini\ldots}

On a donc
\[
H^1_{C_i}(\deleted{\unsure{T}}S,G)\hookrightarrow H^1_{C_j}(S,G)\qquad\text{si }i\geq j
\]
application injective canonique. \uline{Si $G$ est fini}, alors tous les hom sont bijectifs, et $H^1(S,G)$ est l'ensemble des classes (à isom près) des revêtements étales galoisiens de groupe $G$, donc $\simeq\operatorname{Hom}(\pi_1(S,a),G)/\operatorname{int}(G)$. Notre intérêt ici est dans le cas $G$ infini. Notons en passant
\[
H^1_{C_3}\simeq H^1_{C_2}.
\]
